\hypertarget{contributions}{\subsection{Contributions}}
\label{sec:contributions}
To summarize, our contributions are:
\begin{enumerate}
  \item \textbf{Taxonomy.}
  We taxonomize the vast design space of language model evaluation into scenarios and metrics.
  By stating this taxonomy, we can select systematically from this space, which makes explicit both \textit{our priorities} in benchmark design and \textit{the limitations} in the benchmark at present (see \refsec{missing}).
  \item 
  \textbf{Broad coverage.}
  Given our taxonomy, we select and implement \numcorescenarios \core scenarios, for which we comprehensively measure \nummetriccategories metrics (accuracy, calibration, robustness, fairness, bias, toxicity, efficiency).
  We also include \numcomponents targeted evaluations of skills and risks (\eg knowledge, reasoning, disinformation, copyright), introducing \numnewscenarios new scenarios that have not been previously used in mainstream language model evaluation.
  \item \textbf{Evaluation of existing models.}
  We evaluate \nummodels language models under the standardized conditions of our benchmark, ensuring models can now be directly compared across many scenarios and metrics.
  These models vary in terms of their public accessibility: \numpublicmodels are open, \numlimitedmodels are limited-access, and \numclosedmodels are closed.
 \item \textbf{Empirical findings.}
 Our extensive evaluation yields a host of findings (\refsec{experiments}), which in some cases reinforce findings in the literature and in others produce new knowledge about today's language models.
 These results offer guidance for future language model development and ample opportunities for further analysis. 
  \item \textbf{Interactive results and codebase.}
  We provide a public website with all results, underlying model predictions and adaptation details, along an extensible codebase to support the community in taking \benchmarkname further.\footnote{\benchmarkUrl{}}
\end{enumerate}

\paragraph{Acknowledging the prior work this effort builds on.}
To build our holistic evaluation of language models, we directly build on top of many prior works.
While we advocate for evaluating language models in their totality, \ie centralizing many disparate evaluations, we want to be explicit that the underlying works across the AI community \textbf{should be recognized and cited}, as \benchmarkname would not exist in its current form without them.
In particular, if the results of \benchmarkname are used by future work or new models are evaluated on \benchmarkname, they should cite the works that created the many datasets/evaluations that constitute \benchmarkname.\footnote{We take direct inspiration from and follow the precedent set by GLUE and SuperGLUE \citep{wang2019superglue, wang2019glue}; see \url{https://super.gluebenchmark.com/faq}.} 
For this reason, we provide the BibTeX entries for all of these works in the codebase\footnote{\githubUrl{}} and explicitly acknowledge the associated work for every evaluation on the website.\footnote{\benchmarkUrl{}}